{\actuality} Последние десятилетия характеризуются активным внедрением и распространением беспилотных транспортных средств (БТС) в различных сферах человеческой деятельности. Примеры таких сфер -- логистика, промышленность, сельское хозяйство, горное дело, городская инфраструктура. Эксплуатация БТС четвёртого уровня автоматизации по классификации SAE~J3016 вышла за пределы испытательных полигонов, и с расширением области эксплуатации ужесточаются требования к безопасности их движения. Развитие беспилотного транспорта играет важную роль и в достижении национальных целей и задач развития Российской Федерации. Так, <<Транспортные технологии для различных сфер применения (море, земля, воздух), в том числе беспилотные и автономные системы>> входят в перечень критических технологий Российской Федерации. Транспортная стратегия РФ на период до 2030 года включает организацию движения БТС на автомобильной дороге общего пользования. В рамках Стратегии развития автомобильной промышленности Российской Федерации до 2035 года создана инициатива социально-экономического развития <<Беспилотные логистические коридоры>>, предусматривающая проведение мероприятий по внедрению на автомобильных дорогах Государственной компании интеллектуальных транспортных систем, ориентированных в том числе на обеспечение движения высокоавтоматизированных транспортных средств.

Основой безопасного движения БТС служит подсистема восприятия, строящая по данным бортовых датчиков модель окружающего пространства. К такой модели предъявляются два встречных требования. С одной стороны, она должна быть достаточно точной, чтобы обеспечить решение задач использующими её алгоритмами. С другой -- лишённой избыточной детализации: чем сложнее модель, тем больше времени требует работа с ней, а минимизация задержек в системах автопилотирования критически важна, поскольку любая задержка означает, что решение принимается по устаревшему описанию окружения. Поэтому состав сведений, хранимых в модели, определяется тем, какие задачи по ней решаются. В частности, алгоритмам планирования маршрута требуется модель, описывающая расположение зон, движение в которых возможно, невозможно либо о проходимости которых сведения отсутствуют.

Точность описания каждого из этих классов зон одинаково существенна, а ошибки модели опасны в обе стороны: занижение занимаемой препятствием области ведёт к столкновению, завышение -- к ложным остановкам и снижению транспортной работы, а области, ошибочно отнесённые к свободным при фактическом отсутствии наблюдений, лишают алгоритмы планирования возможности действовать консервативно в условиях окклюзии. При этом ни один тип бортовых датчиков не позволяет построить такую модель в одиночку: каждый ограничен собственным принципом работы -- угловым разрешением, дальностью, чувствительностью к погодным условиям либо отсутствием прямого измерения расстояния. Поэтому современные БТС оснащаются наборами разнородных датчиков, а качество модели определяется тем, каким образом их показания комплексируются в единое непротиворечивое описание окружения.

Задача построения модели окружающего пространства, в том числе и модели проходимости окружающего пространства, в робототехнике широко исследована, однако её решение в условиях ограниченного вычислительного бюджета остаётся предметом исследований. Исследование этой задачи нашло отражение в работах М. Павоне, А. Элфеса, Х. Моравека, С. Труна, Б.М. Миллера, Д.А. Юдина и Л.М.Г. Гонсалвеса. Среди исследований задачи комплексирования данных стоит выделить работы Б. Халеги, Д.П. Николаева, Ю.В. Визильтера.

Наибольший объём сведений об окружающей БТС сцене даёт камера: по угловому разрешению она превосходит остальные бортовые датчики и позволяет получить семантическое описание наблюдаемого -- границы проезжей части, классы объектов, дорожную разметку. Однако изображение задаёт лишь направление на объект и не содержит расстояния до него. Поэтому учёт визуальных данных в модели окружающего пространства требует отдельного этапа -- проецирования результатов обработки изображения, детекций объектов и семантической сегментации, из системы координат изображения на плоскость движения БТС. От точности этого проецирования напрямую зависит корректность положения препятствий в модели.

Существующие способы такого проецирования тяготеют к двум крайностям: высокой точности либо низкой вычислительной стоимости. На одном краю -- методы, опирающиеся на известные калибровки камеры и геометрическое предположение о форме поверхности движения. Распространённое допущение о совпадении этой поверхности с плоскостью нарушается как рельефом местности, так и требованиями к конструкции дорожного покрытия. К этому добавляется отклонение фактической внешней калибровки от измеренной, вызванное загрузкой БТС, работой подвески и смещением датчиков при длительной эксплуатации. Вносимая обоими факторами погрешность растёт с удалением от камеры и на тех дальностях, на которых принимается решение об объезде или остановке, составляет единицы метров. На другом краю -- обучаемые преобразования признаков изображения в вид сверху, не использующие явной модели поверхности. Методы этой группы демонстрируют наивысшее качество проецирования, однако требуют производительного  графического ускорителя для запуска нейросетевой модели, принимающей данные от каждого источника данных; размеченных мультимодальных наборов данных для обучения; а также переобучения при любом изменении состава и расположения датчиков, что для мелкосерийных платформ с уникальной сенсорной конфигурацией практически неприемлемо. Предлагаемы в настоящей работе алгоритмы проецирования визуальных данных занимают промежуточное положение между этими крайностями: предположение о совпадении поверхности движения с плоскостью заменяется оценкой её формы в существенно более широком классе поверхностей, но без привлечения обучаемых преобразований в вид сверху.

Второе ограничение имеет тот же источник -- цену одного дополнительного источника данных, -- но проявляется не при проецировании, а при учёте данных в общей модели. Увеличение числа источников данных -- естественный способ сократить слепые зоны БТС и повысить достоверность модели, однако показания каждого из них должны быть сведены в общую модель, а эта операция выполняется над разделяемым ресурсом и потому в существующих реализациях сериализуется. В результате задержка между регистрацией данных датчиком и их отражением в модели растёт с числом источников. Положение усугубляется тем, что для широкоиспользуемых моделей, описывающих пространство в виде двумерного массива -- сеточных представлений, требуемый размер массива определяется тормозным путём и растёт как четвёртая степень рабочей скорости БТС, вследствие чего предельное число одновременно учитываемых источников с ростом скорости сокращается столь же быстро. Известные способы распараллеливания построения сеточных моделей применяются на других уровнях -- внутри одного обновления карты на графическом ускорителе либо между независимыми слоями, сведение которых в общую сетку всё равно выполняется последовательно, -- а неблокирующие структуры данных к разделяемой сетке локальной карты, следующей за БТС и потому требующей операций сдвига и приведения всех элементов к общему моменту времени, не применялись.

Таким образом, качество и актуальность модели, описывающей проходимость окружающего БТС пространства, ограничены двумя факторами: точностью проецирования визуальных данных на плоскость движения без опоры на предположение о плоской поверхности и в пределах вычислительного бюджета бортового вычислителя, а также ростом задержки учёта данных при увеличении числа источников. Оба ограничения имеют общий корень: вычислительная стоимость известных решений растёт с числом учитываемых источников данных, тогда как ресурсы бортового вычислителя разделяются между всеми подсистемами БТС. Разработка вычислительно эффективных алгоритмов комплексирования разнородных данных, а также структуры данных и алгоритма её обновления, допускающих одновременный учёт показаний нескольких источников, представляет собой актуальную научно-техническую задачу.

{\aim} работы является снижение задержки построения и повышение точности модели проходимости окружающего пространства БТС при одновременном учёте данных нескольких разнородных источников в условиях ограниченных вычислительных ресурсов бортового вычислителя.

Для~достижения поставленной цели необходимо было решить следующие {\tasks}:
\begin{enumerate}[beginpenalty=10000] % https://tex.stackexchange.com/a/476052/104425
  \item Провести обзор существующих моделей окружающего пространства БТС и алгоритмов их построения для выбора базовой модели, описывающей проходимость пространства, и алгоритма её построения, которые будут использованы при разработке дальнейших алгоритмов.
  \item Разработать алгоритмы оценки положения объектов в окружающем пространстве БТС, а также областей проезжей
  части, доступных для движения БТС, на основе визуальных и лидарных данных, предназначенные для работы на бортовом вычислителе БТС в режиме реального времени.
  \item Разработать программный комплекс, обеспечивающий построение модели окружающего пространства БТС на основе комплексирования разнородных данных с бортовых датчиков БТС.
  \item Разработать структуру данных для представления модели окружающего пространства и параллельный алгоритм её обновления несколькими источниками данных.
\end{enumerate}


{\novelty}
\begin{enumerate}[beginpenalty=10000] % https://tex.stackexchange.com/a/476052/104425
  \item Предложен алгоритм оценки положения окружающих ТС на плоскости движения БТС по данным монокулярной камеры, отличающийся от известных алгоритмов вписывания в ориентированную ограничивающую рамку тем, что в его основу положена модель L-образной формы видимой границы ТС, ранее применявшаяся лишь к кластерам точек дальномерных датчиков.
  \item Предложен алгоритм оценки положения окружающих ТС на плоскости движения БТС по комплексированию визуальных и лидарных данных, отличающийся от известных тем, что сегментация лидарного облака на точки земли и точки препятствий выполняется без предположения о плоской либо кусочно-плоской поверхности движения -- на основе оценки высоты поверхности регрессией на гауссовских процессах, -- что повышает точность оценки положения во всём диапазоне рассмотренных дальностей.
  \item Предложен алгоритм проекции семантической сегментации проезжей части с изображения на плоскость движения БТС, отличающийся от известных тем, что проецирование выполняется не на фиксированную плоскость, а на оценённую по лидарным данным поверхность движения в проецируемой области, что снижает ошибку описания свободного пространства на строящейся карте проходимости.
  \item Впервые количественно показано, что для класса алгоритмов построения карты проходимости пространства, не использующих обучаемых преобразований в вид сверху и применимых на бортовом вычислителе БТС, учёт визуальных данных с камер повышает качество карты по сравнению с картой, строящейся исключительно по лидарным данным, причём улучшение проявляется в том числе по метрике, учитывающей задачу планирования пути (PFC-MSE).
  \item Предложены структура данных для представления карты проходимости пространства и алгоритм её обновления, отличающиеся от известных реализаций тем, что операции над разделяемой сеткой выполняются без использования примитивов взаимного исключения, что позволяет обновлять карту нескольким источникам данных одновременно и снижает время задержки между моментом регистрации данных датчиком и моментом обновления карты.
\end{enumerate}

{\influence}. Разработанные алгоритмы и комплекс технических средств, обеспечивающие построение модели окружающего пространства БТС,  реализованные в данной работе, внедрены в беспилотные транспортные средства <<EVO.1>> и <<EVO.3>>, разработанные компанией ООО <<ЭвоКарго>>. Внедрение разработанных в данной диссертации методов подтверждено соот­ветствующими актами. На результаты, полученные в работе, получены свидетельства о государственной регистрации программ для ЭВМ:

\begin{enumerate}
	\item №~2025617212 Программное обеспечение <<Система восприятия N1 V3>>.
	\item №~2025669744 Программное обеспечение <<Способ построения карты проходимости в окрестности высокоавтоматизированного беспилотного грузового транспортного средства>>.
\end{enumerate}

На основные результаты работы получен патент №~2853062 <<Способ построения карты проходимости в окрестности высокоавтоматизированного беспилотного грузового транспортного средства>>.

{\methods} В диссертации используются методы вероятностной робототехники, методы байесовской оценки, вычислительной оптимизации, компьютерного зрения, анализ вычислительной сложности, а также численный эксперимент.

{\defpositions}
\begin{enumerate}[beginpenalty=10000] % https://tex.stackexchange.com/a/476052/104425
  \item На записях показаний датчиков реальных БТС предложенный алгоритм вписывания проекции объекта на виде сверху в ориентированную ограничивающую рамку, опирающийся на модель L-образной формы видимой границы ТС, обеспечивает прирост среднего коэффициента Жаккара оценки положения окружающих ТС относительно сравниваемых общеизвестных алгоритмов вписывания, не использующих обучаемые методы и применённых к тем же входным данным: $0,337$ против $0,328$ у лучшего из них (прирост $2,7$~\%).
  \item Среди алгоритмов оценки положения объектов на плоскости движения БТС, применимых на бортовом вычислителе БТС в режиме реального времени, предложенный алгоритм комплексирования визуальных и лидарных данных, не использующий приближение поверхности дороги плоскостью, обеспечивает наибольший средний коэффициент Жаккара для ТС находящихся на расстоянии до 60 м: прирост $10,0$~\% по сравнению с алгоритмом, использующим только лидарные данные ($0,439$ против $0,399$), и прирост $30,3$~\% по сравнению с алгоритмом, использующим только визуальные данные ($0,439$ против $0,337$).
  \item При построении карты проходимости пространства размером $600\times600$ элементов с разрешением $0,1$~м по данным трёх бортовых камер и двух лидаров с ограничением дальности $40$~м предложенный алгоритм проекции семантической сегментации проезжей части на плоскость движения БТС на основе комплексирования визуальных и лидарных данных в сочетании с базовым лидарным алгоритмом повышает качество карты по сравнению с картой, строящейся исключительно по лидарным данным: среднеквадратичная ошибка относительно эталонной карты снижается в среднем с $0,231$ до $0,218$ (на $5,6$~\%), а среднеквадратичная ошибка стоимости поиска пути (PathFinding Cost Mean Squared Error, PFC-MSE) -- с $1019,3$ до $924,4$ (на $9,3$~\%), причём по метрике PFC-MSE улучшение достигается во всех пяти рассмотренных сценариях движения.
  \item Предложенные структура данных для представления карты проходимости пространства и алгоритм, допускающий неблокирующее выполнение операций над картой, уменьшают коэффициент линейной аппроксимации зависимости времени одной итерации цикла обновления карты $8$-ю источника данных от размера карты в $13,9$ раза по сравнению с базовой реализацией, использующей блокировки при выполнении операций.
  \item Разработанный программный комплекс построения карты проходимости пространства позволяет подключать новый источник данных, в том числе датчик нового типа, не изменяя код построения карты. Число одновременно подключённых источников масштабируемо и ограничено только вычислительными возможностями бортового вычислителя.
\end{enumerate}

{\reliability} полученных результатов обеспечивается использованием строгого математического аппарата и методов математической статистики. Помимо этого, достоверность подтверждается результатами вычислительных и натурных экспериментов и внедрением разработанных методов в БТС, находящиеся в эксплуатации.

Экспериментальная оценка выполнена на собственных наборах данных, зарегистрированных с реальных БТС. Публичные наборы данных, применяемые для оценки алгоритмов трёхмерной детекции и сегментации (KITTI, SemanticKITTI, nuScenes), для решаемых в работе задач неприменимы: они не содержат ни конфигурации датчиков, соответствующей рассматриваемым БТС, ни разметки проходимости пространства, необходимой для формирования эталонных карт.

{\probation}
Основные результаты работы докладывались~на международных конференциях:
\begin{itemize}
	\item 16th International Conference on Machine Vision (ICMV~2023, Ереван, Армения);
	\item 40th ECMS International Conference on Modelling and Simulation (ECMS~2026, Гримстад, Норвегия).
\end{itemize}
 

{\contribution} Все результаты диссертации, вынесенные на защиту, получены автором самостоятельно. Постановка задач и обсуждение результатов проводились совместно с научным руководителем.

\ifnumequal{\value{bibliosel}}{0}
{%%% Встроенная реализация с загрузкой файла через движок bibtex8. (При желании, внутри можно использовать обычные ссылки, наподобие `\cite{vakbib1,vakbib2}`).
    {\publications} Основные результаты по теме диссертации изложены
    в~XX~печатных изданиях,
    X из которых изданы в журналах, рекомендованных ВАК,
    X "--- в тезисах докладов.
}%
{%%% Реализация пакетом biblatex через движок biber
    \begin{refsection}[bl-author, bl-registered]
        % Это refsection=1.
        % Процитированные здесь работы:
        %  * подсчитываются, для автоматического составления фразы "Основные результаты ..."
        %  * попадают в авторскую библиографию, при usefootcite==0 и стиле `\insertbiblioauthor` или `\insertbiblioauthorgrouped`
        %  * нумеруются там в зависимости от порядка команд `\printbibliography` в этом разделе.
        %  * при использовании `\insertbiblioauthorgrouped`, порядок команд `\printbibliography` в нём должен быть тем же (см. biblio/biblatex.tex)
        %
        % Невидимый библиографический список для подсчёта количества публикаций:
        \phantom{\printbibliography[heading=nobibheading, section=1, env=countauthorvak,          keyword=biblioauthorvak]%
        \printbibliography[heading=nobibheading, section=1, env=countauthorwos,          keyword=biblioauthorwos]%
        \printbibliography[heading=nobibheading, section=1, env=countauthorscopus,       keyword=biblioauthorscopus]%
        \printbibliography[heading=nobibheading, section=1, env=countauthorconf,         keyword=biblioauthorconf]%
        \printbibliography[heading=nobibheading, section=1, env=countauthorother,        keyword=biblioauthorother]%
        \printbibliography[heading=nobibheading, section=1, env=countregistered,         keyword=biblioregistered]%
        \printbibliography[heading=nobibheading, section=1, env=countauthorpatent,       keyword=biblioauthorpatent]%
        \printbibliography[heading=nobibheading, section=1, env=countauthorprogram,      keyword=biblioauthorprogram]%
        \printbibliography[heading=nobibheading, section=1, env=countauthor,             keyword=biblioauthor]%
        \printbibliography[heading=nobibheading, section=1, env=countauthorvakscopuswos, filter=vakscopuswos]%
        \printbibliography[heading=nobibheading, section=1, env=countauthorscopuswos,    filter=scopuswos]}%
        %
        \nocite{*}%
        %
        {\publications} Основные результаты по теме диссертации изложены в~\arabic{citeauthor}~печатных изданиях,
        \ifnum \value{citeauthorvak}>0\relax
        \arabic{citeauthorvak} из которых изданы в журналах, рекомендованных ВАК%
        \fi%
        \ifnum \value{citeauthorscopuswos}>0%
            , \arabic{citeauthorscopuswos} "--- в~изданиях, индексируемых Web of~Science и Scopus%
        \fi%
        \ifnum \value{citeauthorconf}>0%
            , \arabic{citeauthorconf} "--- в~тезисах докладов.
        \else%
            .
        \fi%
        \ifnum \value{citeregistered}=1%
            \ifnum \value{citeauthorpatent}=1%
                Зарегистрирован \arabic{citeauthorpatent} патент.
            \fi%
            \ifnum \value{citeauthorprogram}=1%
                Зарегистрирована \arabic{citeauthorprogram} программа для ЭВМ.
            \fi%
        \fi%
        \ifnum \value{citeregistered}>1%
            Зарегистрированы\ %
            \ifnum \value{citeauthorpatent}>0%
            \formbytotal{citeauthorpatent}{патент}{}{а}{}%
            \ifnum \value{citeauthorprogram}=0 . \else \ и~\fi%
            \fi%
            \ifnum \value{citeauthorprogram}>0%
            \formbytotal{citeauthorprogram}{программ}{а}{ы}{} для ЭВМ.
            \fi%
        \fi%
        % К публикациям, в которых излагаются основные научные результаты диссертации на соискание учёной
        % степени, в рецензируемых изданиях приравниваются патенты на изобретения, патенты (свидетельства) на
        % полезную модель, патенты на промышленный образец, патенты на селекционные достижения, свидетельства
        % на программу для электронных вычислительных машин, базу данных, топологию интегральных микросхем,
        % зарегистрированные в установленном порядке.(в ред. Постановления Правительства РФ от 21.04.2016 N 335)
    \end{refsection}%
    \begin{refsection}[bl-author, bl-registered]
        % Это refsection=2.
        % Процитированные здесь работы:
        %  * попадают в авторскую библиографию, при usefootcite==0 и стиле `\insertbiblioauthorimportant`.
        %  * ни на что не влияют в противном случае
        \nocite{vakbib2}%vak
        \nocite{patbib1}%patent
        \nocite{progbib1}%program
        \nocite{bib1}%other
        \nocite{confbib1}%conf
    \end{refsection}%
        %
        % Всё, что вне этих двух refsection, это refsection=0,
        %  * для диссертации - это нормальные ссылки, попадающие в обычную библиографию
        %  * для автореферата:
        %     * при usefootcite==0, ссылка корректно сработает только для источника из `external.bib`. Для своих работ --- напечатает "[0]" (и даже Warning не вылезет).
        %     * при usefootcite==1, ссылка сработает нормально. В авторской библиографии будут только процитированные в refsection=0 работы.
}
